\chapter{Emparejamiento y Localización}
\label{cap:localizacion}

Antes de predecir nada hay que responder dos preguntas: qué viaje del horario está haciendo cada autobús del feed, y
en qué punto de ese viaje se encuentra. La respuesta a la segunda es una hora programada, $\Tpos$: la hora a la que el
horario sitúa al autobús en el lugar donde el GPS dice que está.

\section{Emparejamiento con el Horario}
\label{sec:emparejamiento}

\subsection{Emparejamiento Exacto} La referencia $\rho$ del feed es un \texttt{trip\_id} del GTFS. Con la tabla
\texttt{trip\_aliases} (sección «\nameref{sec:alias}») se obtiene directamente el viaje fusionado $J=\mathrm{alias}(\rho)$.
Al construir el emparejador se consultan a la vez los alias de todas las referencias del feed; para cada viaje $J$ se
guarda la \emph{primera} entrada del feed cuyo alias es $J$. En las pruebas con el feed real, el 100\,\% de los
autobuses se emparejó así.

\subsection{Emparejamiento de Respaldo} Si la base de datos no tiene la tabla de alias, o el viaje no aparece en ella,
se buscan las entradas con la misma clave «$\ell|\nu$» que el viaje y se elige la de salida de cabecera más próxima:
\begin{equation}
e^\ast=\argmin_{e\in C}\ \big|d_0^{\mathrm{feed}}(e)-d_1(J)\big|
\quad\text{sujeto a}\quad \big|d_0^{\mathrm{feed}}(e)-d_1(J)\big|\le 180\ \mathrm{s},
\end{equation}
donde $C$ excluye las entradas que ya tienen un emparejamiento exacto con otro viaje. Si hay empate, gana la primera
encontrada; si ninguna cumple la tolerancia (\texttt{MATCH\_TOLERANCE\_SECONDS}), el viaje no tiene datos en vivo.

\section{Proyección de un Punto Sobre un Segmento}
\label{sec:proyeccion}

Toda la localización se basa en una operación: proyectar la posición del autobús sobre un segmento rectilíneo $AB$
(entre dos puntos del trazado, o entre dos paradas).

\subsection{Plano Local} Se usa una proyección equirrectangular centrada en el autobús, que queda en el origen. Un
punto $Q=(\varphi_Q,\lambda_Q)$ pasa a
\begin{equation}
\label{eq:plano}
x_Q=(\lambda_Q-\lambda)\cdot\cos\varphi\cdot 111\,320,\qquad y_Q=(\varphi_Q-\varphi)\cdot 111\,320,
\end{equation}
en metros, con $(\varphi,\lambda)$ la posición del autobús (\texttt{METERS\_PER\_DEGREE}).

\subsection{Proyección} Sean $\mathbf{a}=(x_A,y_A)$ y $\mathbf{b}=(x_B,y_B)$ y $\mathbf{u}=\mathbf{b}-\mathbf{a}$.
\begin{equation}
\label{eq:proyeccion}
f=\begin{cases}\clamp\!\left(-\dfrac{\mathbf{a}\cdot\mathbf{u}}{\|\mathbf{u}\|^2},\,0,\,1\right) & \text{si } \|\mathbf{u}\|^2>0\\[1ex]
0 & \text{si } \|\mathbf{u}\|^2=0\end{cases}
\qquad
d_\perp=\|\mathbf{a}+f\,\mathbf{u}\|,\qquad L=\|\mathbf{u}\|.
\end{equation}
$f\in[0,1]$ es la fracción del segmento hasta el punto más próximo al autobús, $d_\perp$ es la distancia del autobús a
ese punto y $L$ la longitud del segmento.

\begin{proposicion}
El punto $\mathbf{a}+f\,\mathbf{u}$ con $f$ dado por~\eqref{eq:proyeccion} es el punto del segmento más próximo al
origen.
\end{proposicion}
\begin{proof}
Si $\|\mathbf{u}\|=0$ el segmento es un punto. Si no, $g(t)=\|\mathbf{a}+t\mathbf{u}\|^2=\|\mathbf{a}\|^2+2t\,
\mathbf{a}\cdot\mathbf{u}+t^2\|\mathbf{u}\|^2$ es una parábola convexa con mínimo en
$t^\ast=-\mathbf{a}\cdot\mathbf{u}/\|\mathbf{u}\|^2$. Restringida a $[0,1]$, su mínimo es $\clamp(t^\ast,0,1)$, porque
$g$ es decreciente a la izquierda de $t^\ast$ y creciente a la derecha.
\end{proof}

\subsection{Error de la Aproximación} La proyección~\eqref{eq:plano} usa $\cos\varphi$ en la latitud del autobús. En
un punto a $\Delta\varphi$ radianes de distancia norte–sur, el factor correcto sería $\cos(\varphi+\Delta\varphi)
\approx\cos\varphi\,(1-\Delta\varphi\tan\varphi)$. A la latitud de Bilbao ($\varphi\approx43{,}26^\circ$,
$\tan\varphi\approx0{,}94$) y a 3~km ($\Delta\varphi\approx4{,}7\cdot10^{-4}$~rad), el error relativo es del orden
del 0,04\,\%, poco más de un metro. El factor de 111\,320~m por grado sobrestima en torno a un 0,2\,\% la longitud
de un grado de latitud a esa altura, unos 6~m en 3~km. Ambos errores son despreciables frente a los pasillos de
tolerancia (50 y 150~m). Además, la hora programada se obtiene de la \emph{fracción} $f$ y de las distancias del
propio GTFS, no de distancias medidas en el plano, así que un error de escala uniforme no altera $\Tpos$. El de este
plano no es del todo uniforme (frente al elipsoide WGS84, la longitud queda un 0,16\,\% corta y la latitud un
0,2\,\% larga), pero esa diferencia entre los ejes mueve el punto proyectado menos de un 0,4\,\% de la longitud del
segmento: también despreciable.

\begin{figure}[ht]
\centering
\caption{Proyección del Autobús Sobre un Segmento: Fracción $f$ y Distancia Perpendicular $d_\perp$}\label{fig:proyeccion}
\begin{tikzpicture}[scale=1.15]
  \coordinate (A) at (0,0);
  \coordinate (B) at (7,1.4);
  \coordinate (P) at (3.6,2.1);
  \coordinate (H) at ($(A)!(P)!(B)$);
  \draw[very thick,ruta] (A) -- (B);
  \fill[ruta] (A) circle (2.2pt) node[below left] {$A$};
  \fill[ruta] (B) circle (2.2pt) node[below right] {$B$};
  \fill[gps] (P) circle (2.6pt) node[above] {autobús $(0,0)$};
  \draw[gps,dashed,thick] (P) -- (H) node[midway,right] {$d_\perp$};
  \fill[black] (H) circle (1.6pt);
  \draw[decorate,decoration={brace,amplitude=5pt,mirror},thick,gris] ($(A)+(0.05,-0.25)$) -- ($(H)+(0.05,-0.25)$) node[midway,below=4pt,sloped,font=\small] {$f\cdot L$};
  \draw[decorate,decoration={brace,amplitude=5pt,mirror},thick,gris] ($(A)+(0.1,-0.95)$) -- ($(B)+(0.1,-0.95)$) node[midway,below=4pt,sloped,font=\small] {$L$};
\end{tikzpicture}
\end{figure}

\section{Ventana de Búsqueda}
\label{sec:ventana}

El feed indica el orden $n$ de la próxima parada. Esa información acota dónde puede estar el autobús: en algún punto
poco antes de $p_n$. La búsqueda se limita a los tramos
\begin{equation}
\label{eq:ventana}
i_{\min}=\max(1,\,n-3),\qquad i_{\max}=\min(N-1,\,n+2),
\end{equation}
es decir, del tramo $n-3$ al $n+2$: los tres que llevan hasta la próxima parada y los tres que salen de ella o de
las siguientes (\texttt{SEGMENTS\_BEHIND\_NEXT\_STOP} y \texttt{SEGMENTS\_AHEAD\_OF\_NEXT\_STOP}, 3 y 2, son los
desplazamientos respecto a $n$). La ventana cubre las paradas $p_{i_{\min}},\dots,p_{i_{\max}+1}$, de $p_{n-3}$ a
$p_{n+3}$ cuando el viaje las tiene. Mirar
hacia atrás tolera que el feed cambie de próxima parada con algo de retraso; mirar hacia delante, que lo haga con
antelación. Si $n$ no es el orden de ninguna parada del viaje, no hay posición.

\section{Con Trazado}
\label{sec:trazado}

Se usa este método cuando el viaje tiene trazado y se conocen las distancias $s_i$ de todas las paradas de la
ventana. Sean $S_{\mathrm{lo}}=s_{i_{\min}}$ y $S_{\mathrm{hi}}=s_{i_{\max}+1}$.

\begin{enumerate}
  \item \textbf{Candidatos.} Para cada par de puntos consecutivos del trazado $A=\Sigma_j$, $B=\Sigma_{j+1}$ que no
        quedan fuera de la ventana (es decir, salvo si $\sigma_B<S_{\mathrm{lo}}$ o $\sigma_A>S_{\mathrm{hi}}$), se
        proyecta el autobús con~\eqref{eq:proyeccion} y se calcula su posición a lo largo del trazado,
        $s^\ast=\sigma_A+f\,(\sigma_B-\sigma_A)$.
  \item \textbf{Dos grupos.} El candidato es \emph{esperado} si $s^\ast\le s_n+100$~m
        (\texttt{BEYOND\_NEXT\_STOP\_SLACK\_METERS}) y \emph{más allá} en otro caso. En cada grupo se queda el de menor
        $d_\perp$; con empate, el primero.
  \item \textbf{Desempate.} Se elige el mejor esperado salvo que el mejor de más allá esté al menos 20~m más cerca del
        recorrido (\texttt{AMBIGUITY\_MARGIN\_METERS}):
        \[
        \text{elegido}=\begin{cases}\text{más allá}&\text{si existe y } \big(\text{no hay esperado o }d_\perp^{\text{más allá}}+20<d_\perp^{\text{esperado}}\big)\\
        \text{esperado}&\text{en otro caso.}\end{cases}
        \]
  \item \textbf{Pasillo.} La posición es válida si $d_\perp\le50$~m (\texttt{MAX\_DISTANCE\_FROM\_SHAPE\_METERS}).
  \item \textbf{Hora programada.} Se acota $s^\ast\leftarrow\clamp(s^\ast,S_{\mathrm{lo}},S_{\mathrm{hi}})$ y se
        busca el primer tramo $i\in[i_{\min},i_{\max}]$ con $s^\ast\le s_{i+1}$ (o el último de la ventana). Entonces
        \begin{equation}
        \label{eq:tpos}
        f_i=\begin{cases}\clamp\!\Big(\dfrac{s^\ast-s_i}{s_{i+1}-s_i},0,1\Big)&\text{si }s_{i+1}>s_i\\ 0&\text{en otro caso}\end{cases}
        \qquad
        \Tpos=\redondeo\big(d_i+f_i\,(a_{i+1}-d_i)\big).
        \end{equation}
\end{enumerate}

La hora programada se interpola entre la \emph{salida} de la parada de inicio del tramo y la \emph{llegada} a la de
fin: dentro del tramo el autobús está en marcha, y la espera en la parada no forma parte de él. Como los segmentos del
trazado que tocan el borde de la ventana entran como candidatos y $s^\ast$ se acota a $[S_{\mathrm{lo}},S_{\mathrm{hi}}]$,
un GPS algo más allá de un extremo queda en ese extremo: con un trazado que solo tuviera puntos en las paradas, un
autobús 400~m después de $p_{i_{\max}+1}$ y dentro del pasillo toma $\Tpos=a_{i_{\max}+1}$.

Hacen falta dos grupos porque hay recorridos que vuelven sobre la misma calle: una ida y vuelta por el mismo eje, o
una rotonda. Un punto GPS en esa calle tiene dos proyecciones casi igual de cercanas. El grupo esperado prefiere la
que es coherente con la próxima parada que declara el feed. Solo se acepta una proyección más adelantada si es
claramente mejor (20~m), lo que cubre el caso en que el feed tarda en actualizar la próxima parada. Los 100~m de margen
tras la próxima parada absorben que el feed cambia de parada con retraso.

\begin{figure}[ht]
\centering
\caption{Ventana de Búsqueda Alrededor de la Próxima Parada}\label{fig:ventana}
\begin{tikzpicture}[scale=1.0]
  \draw[line width=6pt,gris!25,line cap=round] (0,0) -- (13,0);
  \draw[line width=6pt,ruta,line cap=round] (1.2,0) -- (11.4,0);
  \foreach \x/\l in {1.2/$p_{n-3}$,3.4/$p_{n-2}$,5.6/$p_{n-1}$,7.8/$p_{n}$,9.0/$p_{n+1}$,10.2/$p_{n+2}$,11.4/$p_{n+3}$} {
    \fill[white] (\x,0) circle (5pt); \draw[ruta,thick] (\x,0) circle (5pt);
    \node[below=7pt,font=\small] at (\x,0) {\l};
  }
  \fill[ruta] (7.8,0) circle (5pt);
  \draw[gps,dashed,thick] (7.8,0) -- (7.8,1.3) -- (8.6,1.3) -- (8.6,0);
  \node[gps,font=\footnotesize,above] at (8.2,1.3) {$+100$ m};
  \fill[gps] (6.6,0.55) circle (3pt);
  \draw[gps,thick] (6.6,0.55) -- (6.6,0);
  \node[gps,font=\footnotesize,above] at (6.6,0.6) {GPS};
  \draw[decorate,decoration={brace,amplitude=5pt},thick] (1.2,-0.95) -- (11.4,-0.95) node[midway,below=6pt,font=\small] {ventana: tramos $i_{\min}=n-3$ a $i_{\max}=n+2$, paradas $p_{n-3}$ a $p_{n+3}$};
\end{tikzpicture}
\nota{Se marca el límite de los candidatos «esperados» alrededor de la próxima parada $p_n$.}
\end{figure}

\section{Sin Trazado}
\label{sec:sin-trazado}

Si el viaje no tiene trazado, o falta la distancia de alguna parada de la ventana, se proyecta sobre las rectas entre
paradas. Para cada tramo $i$ de la ventana, en orden, se proyecta el autobús sobre el segmento $p_ip_{i+1}$ (con las
coordenadas de las paradas) y se conserva el de menor $d_\perp$ (con empate, el primero). Para ese tramo se calculan
\begin{equation}
c_i=\max\big(150,\ 0{,}15\,L_i\big),\qquad \Tpos=\redondeo\big(d_i+f\,(a_{i+1}-d_i)\big),
\end{equation}
y la posición es válida si $d_\perp\le c_i$ (\texttt{MAX\_DISTANCE\_FROM\_ROUTE\_METERS},
\texttt{MAX\_DISTANCE\_FROM\_ROUTE\_SEGMENT\_RATIO}). El pasillo es más ancho porque la calle real se aparta de la
recta entre dos paradas, tanto más cuanto más largo es el tramo. Sin trazado no hay dos grupos ni desempate: se queda
el tramo más cercano, y en un recorrido que vuelve sobre sí mismo la única protección es la ventana.

Desde que la construcción mide los metros de las redes que no los publican (sección «\nameref{sec:trazados-gtfs}»), este
caso solo queda para los viajes que no tienen ningún trazado: el resto se sitúa con el pasillo de 50~m.

\section{Condiciones Para Usar el GPS}
\label{sec:condiciones-gps}

La posición de una entrada $e$ en el viaje $J$ existe solo si se cumplen todas estas condiciones:

\begin{enumerate}
  \item $e$ tiene latitud, longitud, hora de la posición $\tgps$ y orden $n$ de la próxima parada;
  \item el dato es reciente: $\tnow-\tgps\le600$~s (\texttt{MAX\_LOCATION\_AGE\_SECONDS});
  \item el viaje tiene al menos dos paradas;
  \item el viaje ha empezado o está a punto: $\tgps\ge d_1-60$~s (\texttt{NOT\_STARTED\_GRACE\_SECONDS}); si no, solo
        cabe la posición en la cabecera de la sección «\nameref{sec:cabecera}»;
  \item $n$ es el orden de una parada del viaje;
  \item la proyección es válida (secciones «\nameref{sec:trazado}» o «\nameref{sec:sin-trazado}»);
  \item el GPS no contradice a la parada registrada (sección «\nameref{sec:desconfianza}»).
\end{enumerate}

Si se cumplen, la \emph{posición} es la cuádrupla $\big(\Tpos,\ \tgps,\ d_\perp,\ \text{método}\big)$ y el
\emph{desfase en la posición} es
\begin{equation}
\label{eq:desfase-pos}
\delta_{\mathrm{pos}}=\tgps-\Tpos.
\end{equation}
Un valor positivo indica que el autobús llega a ese punto más tarde de lo que dice el horario. Como el horario está
interpolado, este desfase \emph{no} se muestra al usuario tal cual (sección «\nameref{sec:retraso-mostrado}»).

\begin{ejemplo}
\label{ej:posicion}
Considérese el viaje sintético del Apéndice~\ref{ap:vectores}: cinco paradas en línea recta hacia el este a lo largo
del paralelo $43{,}26^\circ$, a 0, 400, 1\,000, 2\,500 y 3\,200~m, con horas 8:00:00, 8:01:30, 8:03:30, 8:08:30 y
8:10:30 ($28\,800$, $28\,890$, $29\,010$, $29\,310$ y $29\,430$~s). Un autobús declara como próxima parada $p_3$
($n=3$) y su GPS está a 700~m del origen y $0{,}0001^\circ$ al norte ($\approx11$~m), con $\tgps=28\,950$.

La ventana es $i_{\min}=\max(1,0)=1$, $i_{\max}=\min(4,5)=4$. La mejor proyección cae en el tramo $p_2p_3$, con
$d_\perp=11{,}132$~m $\le50$ y $s^\ast=700$~m. Entonces $f_2=(700-400)/(1000-400)=0{,}5$ y
$\Tpos=\redondeo(28\,890+0{,}5\cdot(29\,010-28\,890))=28\,950$. El desfase en la posición es
$\delta_{\mathrm{pos}}=0$: el autobús está donde el horario dice. Sin trazado se obtiene el mismo $\Tpos$, por
proyección sobre la recta $p_2p_3$.
\end{ejemplo}

\section{En la Cabecera, Antes de Salir}
\label{sec:cabecera}

Un autobús asignado a un viaje que aún no ha salido no cumple la condición~4. Muchos de ellos esperan en la primera
parada: en la captura del 3 de noviembre de 2026 a las 14:47, 36 de los 259 autobuses estaban antes de su hora y 34 de
ellos a menos de 50~m de su primera parada. Sin posición, la parada de cabecera y las siguientes salían solo con el
horario aunque el autobús estuviera allí. Por eso, cuando $\tgps<d_1-60$, la entrada tiene la posición
\begin{equation}
\label{eq:cabecera}
\Tpos=\min(a_1,d_1),\qquad \tgps'=d_1,\qquad\text{método «cabecera»},
\end{equation}
si el autobús está a 150~m o menos de la primera parada (\texttt{MAX\_DISTANCE\_FROM\_ROUTE\_METERS}) y su posición
se tomó como mucho 60~min antes de su salida, $d_1-\tgps\le3\,600$~s (\texttt{MAX\_WAIT\_AT\_ORIGIN\_SECONDS}). Más lejos puede estar acabando otro viaje, y más
de una hora antes la asignación es poco segura. La hora de la posición pasa a ser la de salida, $d_1$: el autobús no
ha recorrido nada y saldrá, como pronto, a su hora.

Con esta posición la fila sale «en directo», pero su llegada no se calcula desde $\tgps'$: es la de la capa 3
(sección «\nameref{sec:cabecera-eta}»). Tampoco cuenta la primera parada como pasada, aunque el feed ya apunte a la
segunda ($n=2$) mientras el autobús espera. La aplicación lo dice como «En la parada», pero solo en la fila de la
primera parada, decidida por su orden en el recorrido y no por su hora: dos paradas seguidas pueden tener la misma hora
publicada, y en un viaje circular la primera y la última son la misma. Las posiciones en la cabecera no entran en las
medidas del aprendizaje, porque no describen ningún recorrido. El 97\,\% de las veces que una entrada se sitúa así, ese
autobús hace después el viaje desde allí (sección «\nameref{sec:fiabilidad-estados}»).

